\scanpage{291}
\mspage{17}
% archive sheet number: 290

ou encore que la flèche $\pi_0^{X,I_{\mathcal H}}$ :
\[
\mathrm{Hom}(X,I_{\mathcal H})\longrightarrow\mathrm{Hom}(\Gamma_{\mathcal H}(X),\Gamma_{\mathcal H}(I_{\mathcal H})\simeq I)
\]
soit bijective.

\uline{Corollaire} 2.\unsure{10}. \inserted{\deleted{(\unsure{avec les notations (2.8.3)})}} Le foncteur $I\unsure{\longleftarrow|\longmapsto}I_{\unsure{W|w}}$ de $\mathrm{Ens}$ dans $\mathrm{Hot}_W$ est pleinement fidèle, i.e. l'hom d'adjonction
\[
\pi_0(I_W)\longrightarrow I
\]
est bijectif. (Ou encore, $\pi_0$ est un foncteur de localisation.)

En effet, cette bijectivité est \deleted{\unsure{totale|visible}} évidente quand on travaille dans $\mathrm{Cat}$, et s'en déduit aussitôt \deleted{\unsure{pour}} dans $\mathrm{Hot}_W$.

\uline{Démonstration de} 2.8. La formule d'adjonction (2.8.2) résulte de la formule correspondante pour $\mathrm{Cat}$, évidente, et du lemme \unsure{2.5|25}. sur les foncteurs adjoints entre catégories de modèles. L'adjonction dans l'autre sens se réduit à la formule (2.8.4).

On a deux flèches canoniques
\diagram{\begin{tikzcd}[ampersand replacement=\&]
\mathrm{Hom}_W(e_W,X)\arrow[r,shift left,"\alpha"] \& \pi_0(X)\arrow[l,shift left,"\beta"],
\end{tikzcd}}
où $\alpha$ est $\pi_0^{e_W,X}$, compte tenu que $\pi_0(e_W)=\pi_0(e_{\mathrm{Cat}})=e_{\mathrm{Ens}}$, et $\mathrm{Hom}(e_{\mathrm{Ens}},\pi_0(X))\simeq\pi_0(X)$.

La flèche $\beta$ est \deleted{\unsure{déduite de}} la flèche \unsure{compo|comp}
\[
\mathrm{Ob}\,X\xrightarrow{\sim}\mathrm{Hom}_{\mathrm{Cat}}(e_{\mathrm{Cat}},X)\xrightarrow{\mathrm{hot}_W}\mathrm{Hom}_W(e_W,X)
\]
